% -------------------------------------------------------------------------------
% Establish page structure & font.
\documentclass[12pt]{report}

\usepackage[total={6.5in, 9in},
	left=1in,
	right=1in,
	top=1in,
	bottom=1in,]{geometry} % Page structure

\usepackage{graphicx} % Required for inserting images
\graphicspath{{../.images/}} % Any additional images I use (BCU logo, etc) are from here.

\usepackage[utf8]{inputenc} % UTF-8 encoding
\usepackage[T1]{fontenc} % T1 font
\usepackage{float}  % Allows for floats to be positioned using [H], which correctly
                    % positions them relative to their location within my LaTeX code.
\usepackage{subcaption}
\usepackage{csquotes}

% -------------------------------------------------------------------------------
% Declare biblatex with custom Harvard BCU styling for referencing.
\usepackage[
    useprefix=true,
    maxcitenames=3,
    maxbibnames=99,
    style=authoryear,
    dashed=false, 
    natbib=true,
    url=false,
    backend=biber
]{biblatex}

\usepackage[british]{babel}

% Additional styling options to ensure Harvard referencing format.
\renewbibmacro*{volume+number+eid}{
    \printfield{volume}
    \setunit*{\addnbspace}
    \printfield{number}
    \setunit{\addcomma\space}
    \printfield{eid}}
\DeclareFieldFormat[article]{number}{\mkbibparens{#1}}

\addbibresource{Report.bib}

% -------------------------------------------------------------------------------
% To prevent "Chapter N" display for each chapter
\usepackage[compact]{titlesec}
\usepackage{wasysym}
\usepackage{import}

\titlespacing*{\chapter}{0pt}{-2cm}{0.5cm}
\titleformat{\chapter}[display]
{\normalfont\bfseries}{}{0pt}{\Huge}

% -------------------------------------------------------------------------------
% Custom macro to make an un-numbered footnote.

\newcommand\blfootnote[1]{
    \begingroup
    \renewcommand\thefootnote{}\footnote{#1}
    \addtocounter{footnote}{-1}
    \endgroup
}

% -------------------------------------------------------------------------------
% Fancy headers; used to show my name, BCU logo and current chapter for the page.
\usepackage{fancyhdr}
\usepackage{calc}
\pagestyle{fancy}

\setlength\headheight{37pt} % Set custom header height to fit the image.

\renewcommand{\chaptermark}[1]{%
    \markboth{#1}{}} % Include chapter name.


% Lewis Higgins - ID 22133848           [BCU LOGO]                [CHAPTER NAME]
\lhead{Lewis Higgins - ID 22133848~~~~~~~~~~~~~~~\includegraphics[width=1.75cm]{BCU}}
\fancyhead[R]{\leftmark}

% ------------------------------------------------------------------------------
% Used to add PDF hyperlinks for figures and the contents page.

\usepackage{hyperref}

\hypersetup{
    colorlinks=true,
    linkcolor=black,
    filecolor=magenta,
    urlcolor=blue,
    citecolor=black,
}

% ------------------------------------------------------------------------------
\usepackage{xcolor} 
\usepackage{colortbl}
\usepackage{longtable}
\usepackage{amssymb}
% ------------------------------------------------------------------------------
\usepackage{tcolorbox}
\newcommand{\para}{\vspace{7pt}\noindent}
% -------------------------------------------------------------------------------

\title{Title TBC - Find a dataset.}
\author{Lewis Higgins - Student ID 22133848}
\date{May 2025}

% -------------------------------------------------------------------------------

\begin{document}


\makeatletter
\begin{titlepage}
    \begin{center}
        \includegraphics[width=0.7\linewidth]{BCU}\\[4ex]
        {\huge \bfseries CMP6228 - Deep Learning Project}\\[2ex]
        {\large \bfseries  \@title}\\[50ex]
        {\@author}\\[2ex]
        {CMP6228 - Deep Neural Networks}\\[2ex]
        {Module Coordinator: Khalid Ismail}\\[2ex]
        {Word count excluding figures, references and appendices: XXXX / 3000}\\[10ex]
    \end{center}
\end{titlepage}
\makeatother
\thispagestyle{empty}
\newpage


% Page counter trick so that the contents page doesn't increment it.
\setcounter{page}{0}

\tableofcontents
\thispagestyle{empty}

\begin{abstract}
    The brief for this assignment mandates an abstract detailing the main findings of the report,
    as seen previously in CMP6202 - AI \& ML.  
\end{abstract}


\chapter*{Introduction}
\addcontentsline{toc}{chapter}{Introduction}
"This section should summarise and highlight the aim
of the report, a brief description of the selected dataset and the achievements
of the report, and finally describe how the report is organised."

\chapter{Problem statement}
"Provide a detailed description of the dataset
and the deep learning problem." 

\para \textbf{It's likely that you can lift most of this from the proposal.}
The feedback from the proposal should be out before this report is due. Therefore, you get 
a slight advantage on this section. Reference literature heavily here, as your use of literature 
is 25\% of this assignment (20\% overall) on its own.

% ! Draft opportunity for this assignment in Week 11 (No exact date, but try for April 11).

\chapter{Proposed method}
"Provide a full description of the proposed
methodology including the pre-processing phase."

% What deep learning methods will you use? Why?
% How will you preprocess your data?

\section{Data ingestion}
Maybe not necessary to be its own section - if it is, it'll have to be brief given 
the heavily limited word count on this assignment.

\section{Data pre-processing}
% Images in full RGB will enormously increase necessary processing time.
% Things like greyscale conversion could go here. Not sure what else there is with images.
\chapter{Experimental results}
"Provide a detailed description of the
parameter settings, evaluation process, and obtained results."

\section{Hyperparameter tuning}

\begin{itemize}
    \item Epochs \begin{itemize}
        \item Number of times the data is sent through the network in training.
    \end{itemize}
    \item Batch size \begin{itemize}
        \item Data isn't all sent into the model at once but rather in partitioned batches 
        where the size is dictated by this parameter.
    \end{itemize}
    \item Learning rate \begin{itemize}
        \item Considered by some to be the most important hyperparameter. Dictates the 
        amount that the model's weights are updated by during training. If too small,
        the model can get stuck. If too big, it may identify a poor solution in one 
        iteration that then has significant influence on further iterations.
        \item As the NUMBER decreases (0.1 -> 0.01), this is actually an INCREASE 
        in learning rate.
    \end{itemize}
    
\end{itemize}


\section{Evaluation}
Appears that some graphs are expected here. \textbf{The raw image files of said plots seem to be 
expected as part of the submission alongside the code.}

% ! You will want to show graphs of your training and validation loss.

Overfitting can be addressed through regularisation. It can also be addressed by adding dropout layers 
to the network, which help to ensure the model learns rather than memorises data. Dropout will ignore some 
neurons while training.

\section{Results}



\chapter*{Conclusion}
\addcontentsline{toc}{chapter}{Conclusion}

\textbf{May also be called the "summary" according to the brief.} 

% ! Your code is part of this submission. I doubt he'll run it (he literally cannot without the dataset)
% ! but he will look at it. Therefore, don't put coloured comments like these in your code, or remove them before submission.
% ? You should also therefore document your code as you did in previous modules. CMP6202 seems like a good reference point.

% ? The word count also excludes the title page, meaning you can subtract 
\addcontentsline{toc}{chapter}{Bibliography}
\printbibliography

\end{document}